% संपूर्ण मराठी ग्रंथातील प्रकरण 35: परिचय
% हा संपादनयोग्य स्रोतखंड आहे. संकलनासाठी संपूर्ण स्रोतसंग्रहातील
% अक्षररूपे, पूर्वभाग, चिन्हव्याख्या आणि आकृती-स्रोत आवश्यक आहेत.
% मूळ: Open Logic Project; मजकूर आणि रूपांतर CC BY 4.0.
% यंत्रानुवाद, दुरुस्ती आणि तपासणी: OpenAI Codex —
% GPT-5.6 Sol आणि GPT-6 Sol, दोन्ही Ultra effort.
% दुरुस्ती-पुनरावलोकन: GPT-6.1 Sol, Ultra effort.
\chapter{परिचय}\label{lam:int:chap}
\begin{editorial}
या प्रकरणात जेरेमी यांनी लॅम्डा कलनावर लिहिलेली मूळ संक्षिप्त टिपणे आहेत.
त्यातील विभाग एकत्र करायचे आहेत; तसेच लॅम्डा-परिभाष्यतेविषयीचा मजकूर
त्या विषयावरील स्वतंत्र, अधिक तपशीलवार प्रकरणातील मजकुरात मिसळायचा आहे.
\end{editorial}








\section{आढावा}\label{lam:int:ovr:sec}

अॅलॉन्झो चर्च यांनी १९३०च्या दशकाच्या सुरुवातीला लॅम्डा कलनाची रचना
रचनाशील तर्कशास्त्राचा आधार म्हणून केली होती; संगणनक्षम फलनांचे
प्रतिमान म्हणून \emph{नव्हे}. पण ते ट्यूरिंग-संगणनक्षम फलने आणि आंशिक
पुनरावर्ती फलने यांसारख्या संगणनक्षमतेच्या इतर व्याख्यांशी समतुल्य
असल्याचे लवकरच दाखवले गेले. सुरुवातीला हे किंचित अनपेक्षित वाटले,
ही गोष्ट हे वैशिष्ट्यीकरण अधिकच रंजक बनवते.

फलनाला स्वतंत्र नाव देण्याऐवजी, त्याची व्याख्या करणाऱ्या सांकेतिक
व्यक्तीकरणाने थेट त्याचा निर्देश करण्याची लॅम्डा संकेतन ही सोयीची
पद्धत आहे. “$f$ हे $f(x) = x + 3$ ने परिभाषित फलन असू द्या,” असे
म्हणण्याऐवजी “$f$ हे $\lambd[x][(x + 3)]$ फलन असू द्या,” असे म्हणता
येते. दुसऱ्या शब्दांत, $\lambd[x][(x+3)]$ हे फलसाधकाच्या मूल्यात तीन मिळवणाऱ्या
फलनाचे केवळ \emph{नाव} आहे. या व्यक्तीकरणातील $x$ हे स्थानधारक आहे: त्याच फलनाचा निर्देश
$\lambd[y][(y + 3)]$ नेही करता येतो. इतर प्राचले असतानाही हे संकेतन
चालते. उदाहरणार्थ, $g(x, y)$ हे दोन चरांचे फलन आणि $k$ ही नैसर्गिक
संख्या असेल, तर $\lambd[x][g(x,k)]$ हे कोणत्याही~$x$ ला
~$g(x, k)$ शी जोडणारे फलन आहे.

सांकेतिक व्यक्तीकरणावरून फलन परिभाषित करण्याच्या या पद्धतीला
\emph{लॅम्डा अमूर्तीकरण} म्हणतात. लॅम्डा अमूर्तीकरणाची दुसरी बाजू
म्हणजे \emph{उपयोजन}: आपल्याकडे एखादे फलन $f$ असेल (समजा, नैसर्गिक
संख्यांवर परिभाषित), तर त्याचे 2 सारख्या कोणत्याही मूल्यावर
\emph{उपयोजन} करता येते. अर्थात, प्रचलित संकेतनात आपण मिळणाऱ्या
मूल्यासाठी~$f(2)$ लिहितो.

लॅम्डा अमूर्तीकरण आणि उपयोजन एकत्र केल्यास काय होते? मग अमूर्तीकृत
चराच्या जागी ज्यावर उपयोजन केले ते मूल्य “बसवून” परिणामी
व्यक्तीकरण सोपे करता येते. उदाहरणार्थ,
\[(\lambd[x][(x + 3)])(2)
\]
हे~$2 + 3$ असे सोपे करता येते.

आतापर्यंत आपण प्रचलित संकल्पनांसाठी नवे संकेतन मांडण्यापलीकडे काही
केलेले नाही. मात्र लॅम्डा कलन संचसैद्धान्तिक दृष्टिकोनापासून अधिक
मूलगामी फारकत घेते. या चौकटीत:
\begin{enumerate}
\item सर्वकाही फलन दर्शवते.
\item लॅम्डा अमूर्तीकरणाने फलने परिभाषित करता येतात.
\item कोणत्याही गोष्टीचे इतर कोणत्याही गोष्टीवर उपयोजन करता येते.
\end{enumerate}
उदाहरणार्थ, $F$ हे लॅम्डा कलनातील पद असेल, तर $F(F)$ अर्थपूर्ण
असल्याचे नेहमी गृहीत धरले जाते. या उदार चौकटीला
\emph{प्रकारविरहित} लॅम्डा कलन म्हणतात; येथे “प्रकारविरहित” याचा अर्थ
“कशाचे कशावर उपयोजन करता येईल यावर बंधन नाही” असा आहे.

\begin{digress}
\emph{प्रकारयुक्त} लॅम्डा कलनही आहे; तो प्रकारविरहित आवृत्तीचा
एक महत्त्वाचा पर्याय आहे. प्रकारयुक्त लॅम्डा कलन अनेक बाबतींत
प्रकारविरहित कलनासारखे असले, तरी त्याची अभिजात संचसैद्धान्तिक
चौकटीशी सांगड घालणे खूप सोपे आहे आणि त्याचे काही गुणधर्मही
खूप वेगळे आहेत.

लॅम्डा कलनावरील संशोधन सैद्धान्तिक संगणकशास्त्रात आणि
प्रोग्रामिंग भाषांच्या रचनेत मध्यवर्ती ठरले आहे. जॉन मॅकार्थी
यांनी १९५०च्या दशकात तयार केलेली LISP ही या कल्पनांचा
प्रभाव असलेल्या भाषेचे सुरुवातीचे उदाहरण आहे.
\end{digress}











\section{लॅम्डा कलनाची विन्यासमीमांसा}\label{lam:int:syn:sec}

सुरुवातीला $x$, $y$, $z$,~\dots अशी चरांची मालिका आणि $a$, $b$,
$c$,~\dots अशी काही स्थिरांक चिन्हे घेतली जातात. पदांचा संच पुढीलप्रमाणे
विगमनाधारित रीतीने परिभाषित केला जातो:
\begin{enumerate}
\item प्रत्येक चर हे पद आहे.
\item प्रत्येक स्थिरांक हे पद आहे.
\item $M$ आणि $N$ ही पदे असतील, तर $(MN)$ हेही पद आहे.
\item $M$ हे पद आणि $x$ हे चर असेल, तर $(\lambd[x][M])$ हे
  पद आहे.
\end{enumerate}
$(MN)$ या रूपातील पदांना \emph{उपयोजने}, तर $(\lambd[x][M])$
या रूपातील पदांना \emph{अमूर्तीकरणे} म्हणतात.

एकही स्थिरांक नसलेल्या प्रणालीला \emph{शुद्ध} लॅम्डा कलन म्हणतात.
आपण मुख्यतः शुद्ध $\lambd$-कलनात काम करणार असल्याने सर्व लहान इंग्रजी
अक्षरे चरे दर्शवतील. $\lambd$-कलनातील पदे दर्शवण्यासाठी आपण मोठी
इंग्रजी अक्षरे ($M$, $N$ इत्यादी) वापरू.

आपण संकेतनाच्या काही रूढी पाळू:

\begin{conv}
\begin{enumerate}
\item कंस वगळले असता उपयोजन डावीकडून उजवीकडे होते. उदाहरणार्थ,
  $M$, $N$, $P$ आणि $Q$ ही पदे असतील, तर $MNPQ$ हे
  $(((MN)P)Q)$ चे संक्षिप्त रूप आहे.
\item पुन्हा, कंस वगळले असता लॅम्डा अमूर्तीकरणाला शक्य तितकी
  व्यापक व्याप्ती द्यायची. उदाहरणार्थ, $\lambd[x][MNP]$ याचा
  अर्थ $(\lambd[x][((MN)P)])$ असा घ्यायचा.
\item एका लॅम्डाने अनेक चरे अमूर्त करता येतात. उदाहरणार्थ,
  $\lambd[xyz][M]$ हे
  $\lambd[x][\lambd[y][\lambd[z][M]]]$ चे संक्षिप्त रूप आहे.
\end{enumerate}
\end{conv}

उदाहरणार्थ,
\[\lambd[xy][xxyx \lambd[z][xz]]
\]
हे पुढीलचे संक्षिप्त रूप आहे:
\[\lambd[x][\lambd[y][((((xx)y)x)(\lambd[z][(xz)]))]].
\]
या रूढी पाठ करा. सुरुवातीला त्या तुमची डोकेदुखी वाढवतील;
पण हळूहळू त्यांची सवय होईल आणि मग कंसांच्या पसाऱ्यापेक्षा
त्या कमी त्रासदायक वाटतील.


इतर चरांचा बंधन-संबंध न बदलता ज्या दोन पदांत फक्त
बद्ध चरांची नावे बदलली आहेत, त्यांना
$\alpha$-सममूल्य म्हणतात; उदाहरणार्थ, $\lambd[x][x]$ आणि
$\lambd[y][y]$. त्यांना “तेच” पद मानणे सोयीचे ठरेल; म्हणजे
$M$ आणि $N$ ही पदे एकच आहेत असे म्हणताना “बद्ध चरांच्या
नामांतराचा फरक सोडून” ती एकच आहेत, असाही अर्थ अभिप्रेत असेल.
हे नामांतर करताना मुक्त चर पकडले जाणार नाहीत याची काळजी घेतली जाते.
एखाद्या चराचा वापर त्याच नावाच्या चराला अमूर्त करणाऱ्या
$\lambd$ च्या व्याप्तीत असेल, तर तो वापर “बद्ध” असतो;
असे बंधन नसलेला वापर “मुक्त” असतो. एकाच वापराला अनेक
अमूर्तीकरणे लागू असतील, तर सर्वांत आतले त्याच नावाचे अमूर्तीकरण
त्याला बांधते. आधीच्या उदाहरणात मुक्त चरे नाहीत; पण
\[(\lambd[z][yz])x
\]
यात $y$ आणि $x$ मुक्त आहेत, तर $z$ बद्ध आहे.











\section{लॅम्डा पदांचे न्यूनीकरण}\label{lam:int:red:sec}


लॅम्डा पदांचे काय करता येते? ती सोपी करता येतात. $M$ आणि $N$ ही
कोणतीही लॅम्डा पदे आणि $x$ हे कोणतेही चर असो. आदेशनानंतर $N$ मधील
मुक्त चरांशी अडथळा निर्माण करतील अशा $M$ मधील बद्ध चरांची आधी
नावे बदलून, मग $M$ मधील~$x$ च्या प्रत्येक मुक्त वापराच्या जागी $N$ आदेशल्यास जो परिणाम
मिळतो, तो आपण $\Subst{M}{N}{x}$ ने दर्शवू. उदाहरणार्थ,
\[\Subst{(\lambd[w][xxw])}{yyz}{x} = \lambd[w][(yyz)(yyz)w].
\]

\begin{digress}
आदेशनासाठी $[N/x]M$, $[x/N]M$ आणि $M[x/N]$ ही पर्यायी
संकेतनेही वापरली जातात. गोंधळ टाळा!
\end{digress}

सहजप्रज्ञेने पाहता, $(\lambd[x][M])N$ आणि $\Subst{M}{N}{x}$
यांचा अर्थ एकच आहे; पहिले पद बदलून दुसरे करण्याच्या क्रियेला
\emph{$\beta$-संकुचन} म्हणतात. $(\lambd[x][M])N$ ला
\emph{रेडेक्स}, आणि $\Subst{M}{N}{x}$ ला त्याचे
\emph{संकुचनफल} म्हणतात. सर्वसाधारणपणे, एखाद्या उपपदाचे
$\beta$-संकुचन करून $P$ हे पद~$P'$ असे बदलता येत असेल, तर
$P$ चे \emph{एका पायरीत $\beta$-न्यूनीकरण होऊन $P'$ मिळते}
असे म्हणतो आणि $P \redone P'$ लिहितो. $P$ पासून अशा
एक-पायरी न्यूनीकरणांच्या काही पायऱ्यांनी (एकही नसेल तरी)
~$P'$ मिळत असेल, तर $P$ चे~$P'$ पर्यंत
\emph{$\beta$-न्यूनीकरण होते}; संकेतांत, $P \red P'$.
ज्या पदाचे आणखी $\beta$-न्यूनीकरण करता येत नाही, त्याला
\emph{$\beta$-अन्यूनीकरणीय} किंवा \emph{$\beta$-सामान्य}
म्हणतात. संदर्भ स्पष्ट असेल तेव्हा “$\beta$-न्यूनीकरण होते”
इत्यादींऐवजी फक्त “न्यूनीकरण होते” असे म्हणू.

काही उदाहरणे पाहू.
\begin{enumerate}
\item आपल्याला मिळते:
\begin{align*}
(\lambd[x][xxy]) \lambd[z][z] & \redone (\lambd[z][z])(\lambd[z][z]) y \\
& \redone (\lambd[z][z]) y \\
& \redone y.
\end{align*}
\item पद “सोपे” केल्याने ते अधिक गुंतागुंतीचे होऊ शकते:
\begin{align*}
(\lambd[x][xxy])(\lambd[x][xxy]) & \redone (\lambd[x][xxy])(\lambd[x][xxy])y \\
& \redone (\lambd[x][xxy])(\lambd[x][xxy])yy \\
& \redone \dots
\end{align*}
\item पद तसेच राहूही शकते:
\[(\lambd[x][xx])(\lambd[x][xx]) \redone (\lambd[x][xx])(\lambd[x][xx]).
\]
\item शिवाय, काही पदांचे एकाहून अधिक प्रकारे न्यूनीकरण करता येते;
  उदाहरणार्थ,
\[(\lambd[x][(\lambd[y][yx]) z)] v \redone (\lambd[y][yv]) z
\]
हे सर्वांत बाहेरच्या उपयोजनाचे संकुचन केल्याने मिळते; आणि
\[(\lambd[x][(\lambd[y][yx]) z)] v \redone (\lambd[x][zx]) v
\]
हे सर्वांत आतल्या उपयोजनाचे संकुचन केल्याने मिळते. मात्र येथे
लक्षात घ्या की दोन्ही पदांचे पुढे न्यूनीकरण होऊन तेच पद, $zv$, मिळते.
\end{enumerate}

शेवटच्या उदाहरणात मिळणारा एकच अंतिम परिणाम हा योगायोग नाही;
तो लॅम्डा कलनाचा “चर्च--रॉसर गुणधर्म” म्हणून ओळखला जाणारा
सखोल आणि महत्त्वाचा गुणधर्म दाखवतो.











\section{चर्च--रॉसर गुणधर्म}\label{lam:int:cr:sec}

\begin{thm}
\label{lam:int:cr:thm:church-rosser}
$M$, $N_1$ आणि $N_2$ ही पदे असून $M \red N_1$ आणि $M \red
N_2$ असे समजा. मग असे $P$ पद आहे की $N_1 \red P$ आणि $N_2 \red P$.
\end{thm}

\begin{cor}
$M$ चे न्यूनीकरण सामान्य रूपापर्यंत करता येते असे समजा. मग हे
सामान्य रूप एकमेव आहे.
\end{cor}

\begin{proof}
$M \red N_1$ आणि $M \red N_2$ असेल, तर मागील प्रमेयामुळे
असे पद~$P$ आहे की $N_1$ आणि $N_2$ या दोन्हींचे न्यूनीकरण
होऊन~$P$ मिळते. $N_1$ आणि~$N_2$ ही दोन्ही सामान्य रूपात
असतील, तर हे केवळ $N_1 \ident P \ident N_2$ असेल तेव्हाच शक्य आहे.
\end{proof}

शेवटी, दोन पदे $M$ आणि~$N$ यांचे न्यूनीकरण होऊन तेच पद मिळत
असेल, तर त्यांना \emph{$\beta$-सममूल्य}, किंवा फक्त
\emph{सममूल्य}, म्हणू; म्हणजेच असे $P$ असेल की $M
\red P$ आणि $N \red P$. हे $M \equal[\beta] N$ असे लिहितात.
\ref{lam:int:cr:thm:church-rosser} वापरून $\equal[\beta]$ हा
समतुल्यता संबंध असल्याचे तपासता येते; शिवाय, कोणत्याही $M$
आणि~$N$ साठी $M \red N$ किंवा $N \red M$ असेल, तर
$M \equal[\beta] N$. (खरे तर, हा गुणधर्म असलेला
$\equal[\beta]$ हा \emph{लहानात लहान} समतुल्यता संबंध आहे,
हे दाखवता येते.)











\section{करीकरण}\label{lam:rep:cur:sec}

$\lambd[x][M]$ हे $\lambd$-अमूर्तीकरण एकच फलसाधक घेणारे
फलन दर्शवते; अनेक फलसाधक घेणारी फलने परिभाषित करायची असतील,
तर ही मोठी मर्यादा ठरते. एक मार्ग म्हणजे जोड्या, तिकड्या इत्यादी
घडवता येतील असा $\lambd$-कलनाचा विस्तार करणे. मग, उदाहरणार्थ,
त्रिस्थानी फलन $\lambd[x][M]$ ला त्याचा फलसाधक म्हणून
तिकडी अपेक्षित असेल. मात्र हे \emph{करीकरणाने} करणे अधिक सोयीचे आहे.

एक उदाहरण पाहू. क्षणभर $\lambd$-कलनात $+$ ही क्रिया आहे
असे मानू. बेरीज हे $2$-स्थानी फलन आहे, म्हणजे ते दोन फलसाधक
घेते. पण $\lambd$-अमूर्तीकरण आपल्याला एकच फलसाधक घेणारी फलने
देते: $\lambd[(x,y)][(x+y)]$ यांसारख्या व्यक्तीकरणांना विन्यासात
परवानगी नाही. तरी $\lambd[y][(x+y)]$ ने दिलेले एकस्थानी
फलन~$f_x(y)$ घेता येते; ते आपल्या एकमेव फलसाधक~$y$ मध्ये
$x$ मिळवते. खरे तर हे एक फलन नाही, तर प्रत्येक संख्या~$x$
साठी “$x$ मिळवा” अशा वेगवेगळ्या फलनांचे कुटुंब आहे. आता
$\lambd[x][f_x]$ असे आणखी एक एकस्थानी फलन~$g$ परिभाषित
करता येते. फलसाधक $x$ ला उपयोजल्यास $g(x)$ हे
फलन~$f_x$ परत करते---म्हणजे $g$ ची मूल्ये इतर फलने आहेत.
आता $g$ चे $x$ ला आणि मिळालेल्या फलनाचे~$y$ ला उपयोजन
केल्यास $(g(x))y = f_x(y) = x+y$ मिळते. अशा प्रकारे हे एकस्थानी फलन द्विस्थानी बेरीज फलनाचेच काम करू शकते.
“करीकरण” म्हणजे द्विस्थानी फलनाचे, ज्याची मूल्ये एकस्थानी
फलने आहेत अशा, एकस्थानी फलनात रूपांतर करण्याची ही युक्ती.


आता $\lambd$-कलनाच्या विन्यासात बसणारे उदाहरण पाहू.
या उदाहरणातील अमूर्तीकरणाचे दुसरे चर आदेशित पहिल्या पदात
मुक्त नसावे; आवश्यक असल्यास त्याचे आधी नामांतर करा.
$f(x,y) = x$ हे फलन कसे दर्शवायचे? दोन फलसाधक घेऊन
पहिला परत करणारे फलन परिभाषित करायचे असेल, तर
$\lambd[x][\lambd[y][x]]$ लिहिता येते. हे अक्षरशः
फलसाधक~$x$ घेऊन फलन~$\lambd[y][x]$ परत करणारे फलन आहे.
$\lambd[y][x]$ हे फलन आणखी एक फलसाधक~$y$ घेते,
पण तो बाजूला ठेवून नेहमी~$x$ परत करते. दोन फलसाधकांना
$\lambd[x][\lambd[y][x]]$ उपयोजल्यास काय होते ते पाहू:
\begin{align*}
  (\lambd[x][\lambd[y][x]])MN
  \bredone & (\lambd[y][M])N \\
  \bredone & M
\end{align*}

पुढील सर्वसाधारण मांडणीत अमूर्तीकरणातील चरे परस्पर भिन्न
आणि उपयोजल्या जाणाऱ्या सर्व पदांतील मुक्त चरांपासून वेगळी
निवडली आहेत; आवश्यक नामांतर आधी केले आहे.
सर्वसाधारणपणे, $x_1$, \dots,~$x_n$ ही प्राचले असलेले
आणि पद~$N$ ने परिभाषित होणारे फलन लिहायचे असेल,
तर $\lambd[x_1][\lambd[x_2][\ldots\lambd[x_n][N]]]$
लिहिता येते. त्याला $n$ फलसाधक उपयोजल्यास मिळते:
\begin{multline*}
  (\lambd[x_1][\lambd[x_2][\ldots\lambd[x_n][N]]]) M_1 \dots M_n \bredone\\
  \begin{aligned}
  \bredone {} & (\Subst{(\lambd[x_2][\ldots\lambd[x_n][N]])}{M_1}{x_1}) M_2
  \dots M_n\\
   \eqs {} & (\lambd[x_2][\ldots\lambd[x_n][\Subst{N}{M_1}{x_1}]]) M_2
            \dots M_n \\
  \vdots & \\
  \bredone {} &\Subst{\Subst{N}{M_1}{x_1}\ldots}{M_n}{x_n}
  \end{aligned}
\end{multline*}
शेवटच्या ओळीचा अक्षरशः अर्थ फलनाच्या व्याख्येच्या देहात
$x_i$ च्या जागी $M_i$ आदेशणे असा आहे; फलनाला अनेक
फलसाधक उपयोजल्यावर आपल्याला नेमके हेच अपेक्षित असते.










\section{लॅम्डा-परिभाष्य अंकगणितीय फलने}\label{lam:int:rep:sec}

लॅम्डा कलन संगणनाचे प्रतिमान कसे ठरू शकते? सुरुवातीला तर या
विधानाचा अर्थ काय, हेही स्पष्ट नाही. नैसर्गिक संख्यांवरील
संगणनक्षमतेची चर्चा करायची असेल, तर त्या संख्यांचे योग्य
निरूपण शोधावे लागेल. आश्चर्यकारकरीत्या चांगले काम करणारे
एक निरूपण पुढे दिले आहे.

\begin{defn}
प्रत्येक नैसर्गिक संख्या~$n$ साठी \emph{चर्च अंक}
$\num{n}$ हे लॅम्डा पद $\lambd[x][\lambd[y][(x(x(x(\dots
x(y)))))]]$ असे परिभाषित करा; अमूर्तीकरणाच्या देहात एकूण $n$ वेळा $x$ येते.
\end{defn}

$\num{n}$ ही पदे “पुनरावर्तक” आहेत: $f$ हा निविष्ट दिल्यास
$\num{n}$ हे $y$ ला $f^n(y)$ शी जोडणारे फलन परत करते.
प्रत्येक अंक सामान्य रूपात आहे, हे लक्षात घ्या. आता
एखादे लॅम्डा पद नैसर्गिक संख्यांवरील फलनाचे “संगणन करते”
याचा अर्थ सांगता येईल.

\begin{defn}
$f(x_0, \dots, x_{k-1})$ हे $\Nat$ पासून $\Nat$ मध्ये
जाणारे $k$-स्थानी आंशिक फलन असो. एखादे बंद $\lambd$-पद~$F$
हे~$f$ ला \emph{लॅम्डा-परिभाषित करते} असे म्हणतो, जर आणि फक्त जर
नैसर्गिक संख्यांच्या प्रत्येक $n_0$, \dots,~$n_{k-1}$
अनुक्रमासाठी पुढील दोन्ही अटी पूर्ण होत असतील: $f(n_0, \dots, n_{k-1})$ परिभाषित असेल तर
\[F\, \num{n_0}\, \num{n_1} \dots \num{n_{k-1}} \red \num{f(n_0, n_1, \dots,
  n_{k-1})}
\]
असे असते; आणि तसे नसेल तर $F\, \num{n_0}\, \num{n_1}
\dots \num{n_{k-1}}$ याला सामान्य रूप नसते.
\end{defn}

\begin{thm}
\label{lam:int:rep:thm:lambda-def}
फलन $f$ हे आंशिक संगणनक्षम फलन असते, जर आणि फक्त जर
एखाद्या लॅम्डा पदाने ते लॅम्डा-परिभाषित असते.
\end{thm}

\begin{explain}
हे प्रमेय काहीसे थक्क करणारे आहे. संगणनाचे प्रतिमान म्हणून
लॅम्डा कलन अगदी साधे कलन आहे; त्यात लॅम्डा अमूर्तीकरण आणि
उपयोजन एवढ्याच क्रिया आहेत!{} तरी या मोजक्या साधनांतून
कोणतीही संगणनप्रक्रिया साकार करता येते.
\end{explain}











\section{लॅम्डा-परिभाष्य फलने संगणनक्षम असतात}\label{lam:int:cmp:sec}

\begin{thm}
\label{lam:int:cmp:thm:lambda-computable}
आंशिक फलन $f$ हे एखाद्या लॅम्डा पदाने लॅम्डा-परिभाषित
असेल, तर ते संगणनक्षम आहे.
\end{thm}

\begin{proof}
फलन~$f$ हे लॅम्डा पद~$X$ ने लॅम्डा-परिभाषित आहे असे समजा.
~$f$ चे संगणन करण्याची एक अनौपचारिक पद्धत सांगू.
$m_0$, \dots,~$m_{n-1}$ ही निविष्ट मिळाल्यावर
$X \num m_0 \ldots \num m_{n-1}$ हे पद लिहा आणि आधी तेच चर्च अंक आहे का ते तपासा.
मग मूळ पदाची सर्व एका-पायरीतील न्यूनीकरणे लिहून वृक्ष तयार करा;
त्याखाली त्या प्रत्येक पदाची सर्व एका-पायरीतील न्यूनीकरणे
(म्हणजे मूळ पदाची दोन-पायरीतील न्यूनीकरणे) लिहा; आणि
प्रत्येक पातळी पूर्ण करून असेच पुढे चालू ठेवा. पदातील redex
जागा सांत असतात; प्रत्येक पायरीतील आवश्यक नामांतरासाठी
ठरलेली ताजी नावे घेतल्याने प्रत्येक पातळीही सांत राहते.
नामांतरापर्यंत कोणताही चर्च अंक मिळाला, तर संबंधित संख्या
उत्तर म्हणून परत करा. फलनाचे मूल्य अपरिभाषित असेल, तर
असा अंक मिळत नाही आणि शोध अखंड चालू राहतो.

चर्च यांच्या प्रबंधाचा आधार घेतल्यास हे फलन संगणनक्षम
असल्याचे कळते. प्रमेय सिद्ध करण्याचा अधिक चांगला मार्ग
म्हणजे या शोधपद्धतीचे पुनरावर्ती वर्णन देणे. उदाहरणार्थ,
“$\fn{IsASubterm}$,” “$\fn{Substitute}$,”
“$\fn{ReducesToInOneStep}$,” “$\fn{ReductionSequence}$,”
“$\fn{Numeral}$” इत्यादी आदिम पुनरावर्ती फलने आणि
संबंधांची मालिका परिभाषित करता येईल. मग
$f(m_0, \dots, m_{n-1})$ चे संगणन करणारी आंशिक
पुनरावर्ती पद्धत म्हणजे $X \num{m_0} \dots \num{m_{n-1}}$
पासून सुरू होऊन अंकापाशी संपणारा एका-पायरीतील
न्यूनीकरणांचा अनुक्रम शोधणे आणि त्या अंकाशी संबंधित
संख्या परत करणे. तपशील मोठे आणि कंटाळवाणे आहेत,
पण इतरथा नेहमीचे आहेत.
\end{proof}











\section{संगणनक्षम फलने लॅम्डा-परिभाष्य असतात}\label{lam:int:lrp:sec}

\begin{thm}
\label{lam:int:lrp:thm:computable-lambda}
प्रत्येक संगणनक्षम आंशिक फलन लॅम्डा-परिभाष्य असते.
\end{thm}

\begin{proof}

प्रत्येक आंशिक संगणनक्षम फलन~$f$ साठी बंद प्रतिनिधी~$F$
शोधायचा आहे. क्लिनीच्या प्रमाण रूप प्रमेयानुसार, total आदिम
पुनरावर्ती अंकफलने दर्शवणे, total आदिम पुनरावर्ती zero-test
चे अपरिबद्ध लघुतमीकरण करणे, आणि त्यानंतर योग्य output
काढणे पुरेसे आहे. Total आदिम पुनरावर्ती फलनांसाठी आरंभीची
फलने, total संयोजन आणि total आदिम पुनरावर्तन यांच्या पुढील
constructions वापरू. Search चे मूल्य अपरिभाषित असताना
output extractor ते टाकून देऊ नये म्हणून शेवटचे उपयोजन
ते search आधी force करेल. ही शेवटची construction
लघुतमीकरणाच्या विभागात देऊन या प्रमेयाची सिद्धता पूर्ण करू.
सर्व आंशिक फलनांसाठी सरळ संयोजन किंवा सरळ recursion
योग्य असल्याचे आधीच गृहीत धरण्याची गरज नाही.
\end{proof}

उर्वरित पुरावा अधिक वाचनीय व्हावा म्हणून आपण अधिक प्रचलित
संकेतन वापरू. उदाहरणार्थ, $Mxyz$ ऐवजी $M(x, y, z)$ लिहू.
यातून काही अर्थ सूचित होत असला, तरी प्रकारविरहित लॅम्डा
कलनातील पदांना निश्चित स्थानसंख्या नसते, हे लक्षात ठेवा;
म्हणून त्याच पदासाठी~$M$ आपण $M(x,y)$ आणि
$M(x,y,z,w)$ असे दोन्ही लिहू शकतो. पण या संकेतनाने
आपण $M$ ला तीन चरांचे फलन मानत आहोत, हे सूचित होते
आणि पुढील व्याख्यांमागील हेतू अधिक स्पष्ट होतो.
तसेच “$M$ ची व्याख्या $M =
\lambd[x][\lambd[y][\lambd[z][\dots]]]$ ने करा”
याऐवजी “$M$ ची व्याख्या $M(x,y,z) = \dots$ ने करा”
असे म्हणू.











\section{मूलभूत आदिम पुनरावर्ती फलने लॅम्डा-परिभाष्य असतात}\label{lam:int:bas:sec}

\begin{lem}
$\Zero$, $\Succ$ आणि $\Proj{n}{i}$ ही फलने
लॅम्डा-परिभाष्य आहेत.
\end{lem}

\begin{proof}

येथे $\Zero$ हे आधी परिभाषित केलेले एकस्थानी शून्य फलन आहे.
त्याला दर्शवणारे पद
$\lambd[w][\lambd[x][\lambd[y][y]]]$
असे घ्या. बाहेरचे अमूर्तीकरण निविष्ट बाजूला ठेवते आणि
चर्चचा शून्य अंक परत करते.

उत्तरवर्ती फलन~$\Succ$ ची
$\fn{Succ}(u) = \lambd[x][\lambd[y][x(uxy)]]$ अशी
व्याख्या करतात. हे कसे चालते याचा विचार करा: पुनरावर्तक
मानलेल्या प्रत्येक अंकासाठी $\num{n}$ आणि प्रत्येक
फलनासाठी~$f$, $\fn{Succ}(\num{n},f)$ हे असे फलन आहे
की निविष्ट~$y$ वर ते~$y$ पासून सुरू करून $f$ चे $n$
वेळा उपयोजन करते आणि त्यानंतर आणखी एकदा उपयोजन करते.

प्रक्षेपणांबद्दल अधिक काही सांगण्यासारखे नाही:
$\fn{Proj}^n_i(x_0, \dots,
x_{n-1}) = x_i$. दुसऱ्या शब्दांत, आपल्या रूढींनुसार,
$\fn{Proj}^n_i$ हे लॅम्डा पद
$\lambd[x_0][\dots \lambd[x_{n-1}][x_i]]$ आहे.
\end{proof}











\section{लॅम्डा-परिभाष्य फलने संयोजनाच्या अंतर्गत बंद असतात}\label{lam:int:com:sec}

\begin{lem}
लॅम्डा-परिभाष्य फलनांचा वर्ग संयोजनाच्या अंतर्गत बंद असतो.
\end{lem}

\begin{proof}

प्रथम या सरळ मांडणीत सर्व घटक फलने सर्वत्र परिभाषित आहेत
असे घ्या. $f$ हे $h$, $g_0,$ \dots,~$g_{k-1}$ यांच्या संयोजनाने
परिभाषित आहे असे समजा. $h$, $g_0$, \dots,~$g_{k-1}$
ही फलने अनुक्रमे $H$, $G_0$, \dots,~$G_{k-1}$
यांनी लॅम्डा-परिभाषित आहेत असे गृहीत धरल्यास,
~$f$ ला लॅम्डा-परिभाषित करते असे पद~$F$ शोधायचे आहे.
पण~$F$ ची व्याख्या सरळ अशी करता येते:
\[F(x_0, \dots, x_{l-1}) = H(G_0(x_0, \dots, x_{l-1}),
\dots, G_{k-1}(x_0, \dots, x_{l-1})).
\]
प्रत्येक घटक निविष्ट अंकांवर सामान्य रूपातील अंक देतो;
ते सर्व न्यूनीकरण करून मग बाह्य फलनाचे न्यूनीकरण केल्याने
संयोजनाचा योग्य अंक मिळतो. त्यामुळे total फलनांसाठी ही
मांडणी योग्य आहे.
आंशिक फलनांसाठी हा सरळ पुरावा पुरेसा नाही. बाह्य फलन
निविष्ट वापरत नसेल, तर ते अपरिभाषित अंतर्गत संगणनही
टाकून देऊ शकते. म्हणून सामान्य strict partial closure साठी
पुढे पूर्ण होणारे \ref{lam:int:lrp:thm:computable-lambda} वापरतो:
\ref{lam:int:cmp:thm:lambda-computable} ने सर्व घटक आंशिक संगणनक्षम
आहेत; त्यांचे strict संयोजन आंशिक संगणनक्षम आहे; आणि
पूर्ण झालेले पहिले प्रमेय त्याला लॅम्डा प्रतिनिधी देते.
त्या प्रमेयाची पुढील सिद्धता येथे फक्त total closure वापरते,
म्हणून यात चक्रीय अवलंबन नाही.
\end{proof}











\section{लॅम्डा-परिभाष्य फलने आदिम पुनरावर्तनाच्या अंतर्गत बंद असतात}\label{lam:int:pr:sec}


आदिम पुनरावर्तनासाठी आपण टप्प्याटप्प्याने construction करू.
या मुख्य सिद्धतेत $g$ आणि~$h$ ही सर्वत्र परिभाषित फलने असून
त्यांना अनुक्रमे लॅम्डा-परिभाषित करणारी बंद पदे $G'$ आणि $H'$
आधीच दिलेली आहेत. पुढील फलन~$f$ ला दर्शवणारे पद $F'$ हवे:
\begin{align*}
f(0, \vec z) & = g(\vec z) \\
f(x+1, \vec z) & = h(x, f(x,\vec z), \vec z).
\end{align*}
प्रत्येक नैसर्गिक संख्या~$n$ साठी पुढील अटी पुरेशा आहेत:
\begin{align*}
F'(\num{0}, \vec z) & \equiv G'(\vec z) \\
F'(\overline{n+1}, \vec z) & \equiv H'(\num{n}, F'(\num{n}, \vec z), \vec z).
\end{align*}
कारण निविष्ट प्राचलांच्या जागी अंक घातल्यावर विगमनाने
प्रत्येक पायरीवर योग्य सामान्य रूपातील अंक मिळतो.

प्राचले फलनमूल्यात सामावण्यासाठी पुढील अटी पूर्ण करणारे
पद~$F$ शोधणे पुरेसे आहे:
\begin{align*}
F(\num 0) & \equiv G \\
F(\overline {n+1}) & \equiv H(\num{n},F(\num{n}))
\intertext{प्रत्येक नैसर्गिक संख्या $n$ साठी; येथे}
G & = \lambd[\vec z][G'(\vec z)] \text{ आणि}\\
H(u,v) & = \lambd[\vec z][H'(u,v(\vec z),\vec z)].
\end{align*}
आता $F'(u,\vec z)=F(u)(\vec z)$ असे घेतल्यास वरची
प्राचलांसहित recurrence मिळते. येथे मागील value function
आधीच योग्य recursion index साठी घेतलेले आहे; त्याला
फक्त उरलेली प्राचले उपयोजायची आहेत.

पद $F$ परिभाषित करण्यापूर्वी क्रमित जोड्या हाताळण्याची
पद्धत हवी. पुढील उपपत्ती ती देते.

\begin{lem}
असे लॅम्डा पद $D$ आहे की लॅम्डा पदांच्या प्रत्येक जोडी
$M$ आणि~$N$ साठी $D(M,N)(\num{0}) \red M$ आणि
$D(M,N)(\num{1}) \red N$.
\end{lem}

\begin{proof}
आधी लॅम्डा पद $K$ ची पुढीलप्रमाणे व्याख्या करा:
\[K(y) = \lambd[x][y].
\]
म्हणजे $K$ हे पद $\lambd[y][\lambd[x][y]]$ आहे.
दुसऱ्या प्रकारे पाहिले, तर प्रत्येक $M$ साठी $K(M)$
हे कोणतीही निविष्ट घेतल्यावर~$M$ परत करणारे स्थिर फलन आहे.

आता $D(x,y,z)$ ची व्याख्या $D(x,y,z) = z (K(y))x$
अशी करा. मग आपल्याला मिळते:
\begin{align*}
D(M,N,\num 0) & \red \num 0 (K(N)) M \red M \text{ आणि}\\
D(M,N,\num 1) & \red \num 1 (K(N)) M \red K(N) M \red N,
\end{align*}
हेच अपेक्षित होते.
\end{proof}

येथे कल्पना अशी आहे की $D(M,N)$ ही $\tuple{M, N}$
जोडी दर्शवते. $P$ अशी जोडी दर्शवते असे गृहीत धरल्यास,
$P(\num 0)$ आणि $P(\num 1)$ हे अनुक्रमे डावे आणि उजवे
प्रक्षेपण, $(P)_0$ आणि $(P)_1$, दर्शवतात. पुढे आपण
हीच संकेतने वापरू.

\begin{lem}
लॅम्डा-परिभाष्य फलनांचा वर्ग आदिम पुनरावर्तनाच्या
अंतर्गत बंद असतो.
\end{lem}

\begin{proof}
कोणतीही पदे $G$ आणि $H$ दिली असता प्रत्येक नैसर्गिक
संख्या~$n$ साठी पुढील अटी पूर्ण करणारे पद $F$
शोधता येते, हे दाखवायचे आहे:
\begin{align*}
F(\num{0}) & \equiv G \\
F(\num{n+1}) & \equiv H(\num{n}, F(\num{n}))
\end{align*}
साधारण कल्पना म्हणजे अंकांचा पुनरावर्तक म्हणून वापर करून
\emph{जोड्यांचा} अनुक्रम संगणित करणे:
\[\tuple{\num 0, F(\num 0)}, \tuple{\num 1, F(\num 1)}, \dots,
\]
पहिली जोडी म्हणजे $\tuple{\num 0, G}$, हे लक्षात घ्या.
$\tuple{\num{n}, F(\num{n})}$ ही जोडी दिल्यास पुढची जोडी
$\tuple{\num{n+1}, F(\num{n+1})}$ ही
$\tuple{\num{n+1}, H(\num{n}, F(\num{n}))}$ शी
सममूल्य असायला हवी. हे एका पायरीतील संक्रमण करणारे
लॅम्डा पद~$T$ आपण तयार करू.

तपशील असे आहेत. येथे S हे आधी तयार केलेले उत्तरवर्ती
फलन दर्शवते. $T(u)$ ची व्याख्या पुढीलप्रमाणे करा:
\[T(u) = \tuple{S((u)_0), H((u)_0,(u)_1)}.
\]
मग कोणत्याही संख्या~$n$ साठी पुढील विधान सहज तपासता येते:
\[T(\tuple{\num{n}, M}) \red \tuple{\num{n+1}, H(\num{n}, M)}.
\]
वर सुचवल्याप्रमाणे, $G$ आणि $H$ दिली असता
~$F(u)$ ची व्याख्या पुढीलप्रमाणे करा:
\[F(u) = (u(T,\tuple{\num 0, G}))_1.
\]
म्हणजे निविष्ट~$\num{n}$ वर $F$ हे
$\tuple{\num 0, G}$ पासून सुरू करून $T$ चे $n$
वेळा पुनरावर्तन करते आणि मग दुसरा घटक परत करते.
सुरुवातीला आपल्याला मिळते:
\begin{enumerate}
\item $\num{0} (T, \tuple{\num 0, G}) \equiv \tuple{\num{0}, G}$
\item $F(\num{0}) \equiv G$
\end{enumerate}
~$n$ वरील विगमनाने प्रत्येक नैसर्गिक संख्येसाठी
पुढील गोष्टी दाखवता येतात:
\begin{enumerate}
\item $\num{n+1}(T, \tuple{\num{0}, G}) \equiv \tuple{\num{n+1}, F(\num{n+1})}$
\item $F(\num{n+1}) \equiv H(\num{n}, F(\num{n}))$
\end{enumerate}
दुसऱ्या विधानासाठी,
\begin{align*}
F(\num{n+1}) & \red (\num{n+1}(T, \tuple{\num 0, G}))_1 \\
& \equiv (T(\num{n} (T, \tuple{\num 0, G})))_1 \\
& \equiv (T(\tuple{\num{n}, F(\num{n})}))_1 \\
& \equiv (\tuple{\num{n+1}, H(\num{n}, F(\num{n}))})_1 \\
& \equiv H(\num{n}, F(\num{n})).
\end{align*}
येथे तिसऱ्या ओळीत विगमन गृहीतक वापरले आहे.
पहिल्या विधानासाठी,
\begin{align*}
\num{n+1} (T, \tuple{\num 0, G}) &
\equiv T(\num{n} (T, \tuple{\num 0, G})) \\
& \equiv T( \tuple{\num{n}, F(\num{n})}) \\
& \equiv \tuple{\num{n+1}, H(\num{n}, F(\num{n}))} \\
& \equiv \tuple{\num{n+1}, F(\num{n+1})}.
\end{align*}
येथे शेवटच्या ओळीत दुसरे विधान वापरले आहे.
अशा प्रकारे $F(\num 0) \equiv G$ आणि प्रत्येक $n$
साठी $F(\num {n+1}) \equiv H(\num
n, F(\num{n}))$ दाखवले; आपल्याला नेमके हेच हवे होते.
वरील beta-recursion equations कोणत्याही दिलेल्या पदांसाठी
योग्य आहेत. Numeric फलनाच्या प्रतिनिधित्वाचा येथे वापर
total मूळ व step फलनांसाठी केला आहे. सामान्य आंशिक
strict पुनरावर्तनाचा निष्कर्ष \ref{lam:int:cmp:thm:lambda-computable}
आणि पुढे स्वतंत्रपणे पूर्ण होणाऱ्या
\ref{lam:int:lrp:thm:computable-lambda} मधून मिळतो: आंशिक
संगणनक्षम फलने strict आदिम पुनरावर्तनाच्या अंतर्गत बंद आहेत.
Step फलन आधीचे अपरिभाषित value टाकून देऊ शकते,
म्हणून arbitrary पदांसाठीची equations स्वतःहून तो
आंशिक निष्कर्ष सिद्ध करतात असा दावा नाही.
\end{proof}











\section{स्थिरबिंदू संयोजक}\label{lam:int:fix:sec}

तुमच्याकडे लॅम्डा पद $g$ आहे आणि दुसरे पद $k$ हवे
आहे असे समजा; त्यात $k$ हे $gk$ शी
$\beta$-सममूल्य असावे.

आपल्या संकेतनरूढींनुसार पुढील पदे परिभाषित करा;
येथील अमूर्त केलेले चर दिलेल्या पदात मुक्त नसावे.
आवश्यक असल्यास त्याचे आधी नामांतर करा:
\[\fn{diag}(x) = xx
\]
आणि
\[l(x) = g(\fn{diag}(x))
\]
म्हणजे $l$ हे पद $\lambd[x][g(xx)]$ आहे.
$k$ हे पद $ll$ असू द्या. मग आपल्याला मिळते:
\begin{align*}
k & = (\lambd[x][g(xx)])(\lambd[x][g(xx)]) \\
& \red  g((\lambd[x][g(xx)])(\lambd[x][g(xx)])) \\
& = gk.
\end{align*}
आता
\[Y = \lambd[g][((\lambd[x][g(xx)])(\lambd[x][g(xx)]))]
\]
असे घेतल्यास $Yg$ आणि $g(Yg)$ यांचे न्यूनीकरण
होऊन एकच पद मिळते; म्हणून $Yg \equiv_\beta
g(Yg)$. याला “करीचा संयोजक” म्हणतात. त्याऐवजी
\[Y = (\lambd[xg][g(xxg)])(\lambd[xg][g(xxg)])
\]
असे घेतल्यास प्रत्यक्षात $Yg$ चे $g(Yg)$ पर्यंत
न्यूनीकरण होते; हे अधिक प्रबळ विधान आहे.
$Y$ च्या या दुसऱ्या रूपाला “ट्यूरिंगचा संयोजक”
म्हणतात.











\section{लॅम्डा-परिभाष्य फलने लघुतमीकरणाच्या अंतर्गत बंद असतात}\label{lam:int:min:sec}

\begin{lem}
$f(x,y)$ हे लॅम्डा-परिभाष्य आहे असे समजा.
$g$ ची पुढीलप्रमाणे व्याख्या करा:
\[g(x) \simeq \umin{y}{f(x,y)}.
\]
मग $g$ हे लॅम्डा-परिभाष्य आहे.
\end{lem}

\begin{proof}

प्रथम $f$ सर्वत्र परिभाषित आहे असे घ्या. मुख्य कल्पना साधारण अशी आहे. $x$ दिले असता,
~$n$ पासून सुरू करून~$y$ शोधणारे फलन $h_x(n)$
परिभाषित करण्यासाठी स्थिरबिंदू लॅम्डा पद $Y$ वापरू;
मग $g(x)$ म्हणजे $h_x(0)$. फलन~$h_x$ हे
स्थिरबिंदू समीकरणाचे उत्तर म्हणून मांडता येते:
\[h_x(n) \simeq
\begin{cases}
n & \text{जर $f(x,n) = 0$ असेल} \\
h_x(n+1) & \text{जर $f(x,n)>0$ असेल.}
\end{cases}
\]

या total प्रकरणाचे तपशील असे आहेत. $f$ हे लॅम्डा-परिभाष्य असल्याने
कोणते तरी बंद पद $F$ त्याला लॅम्डा-परिभाषित करते.
$D(M, N, \bar{0}) \red M$ आणि $D(M, N, \bar{1})
\red N$ असे असणारे लॅम्डा पद~$D$ आपल्याकडे आहे,
हे लक्षात ठेवा. क्षणभर $x$ स्थिर धरू. $h_x$ ला
लॅम्डा-परिभाषित करण्यासाठी~$x$ वर अवलंबून असणारे
आणि पुढील अट पूर्ण करणारे पद~$H$ शोधायचे आहे:
\[H(\num{n}) \equiv D(\num{n}, H(S(\num{n})), F(x, \num{n})).
\]
स्थिरबिंदू पद~$Y$ वापरून हे करता येते. प्रथम $U$
हे पुढील पद असू द्या:
\[\lambd[h][\lambd[z][D(z,(h(Sz)),F(x,z))]],
\]
आणि मग $H$ हे पद~$YU$ असू द्या. $H$ मधील
एकमेव मुक्त चर~$x$ आहे, हे लक्षात घ्या. वरचे
समीकरण $H$ पूर्ण करते, हे दाखवू.

$Y$ च्या व्याख्येनुसार,
\[H = YU \equiv U(YU) = U(H).
\]
विशेषतः, प्रत्येक नैसर्गिक संख्या~$n$ साठी,
\begin{align*}
H(\num{n}) & \equiv U(H, \num{n}) \\
& \red D(\num{n}, H(S(\num{n})), F(x, \num{n})),
\end{align*}
अपेक्षेप्रमाणे. शेवटच्या ओळीत~$x$ च्या जागी
$\num{m}$ हा अंक आदेशल्यास,
$F(\num{m}, \num{n})$ चे न्यूनीकरण होऊन
$\num{0}$ मिळते तेव्हा हे व्यक्तीकरण $\num{n}$
पर्यंत न्यूनीत होते; आणि $F(\num{m}, \num{n})$
चे न्यूनीकरण होऊन इतर कोणताही अंक मिळतो तेव्हा
ते $H(S(\num{n}))$ पर्यंत न्यूनीत होते.

पुरावा पूर्ण करण्यासाठी $G$ हे
$\lambd[x][H(\num 0)]$ असू द्या. मग $G$
हे~$g$ ला लॅम्डा-परिभाषित करते; म्हणजेच प्रत्येक~$m$
साठी $g(m)$ परिभाषित असेल तर $G(\num m)$ चे
न्यूनीकरण होऊन~$\overline {g(m)}$ मिळते;
Zero witness मिळाल्यास आधीच्या सांत positive tests नंतर
तोच अंक मिळतो. Witness नसल्यास प्रत्येक total test
positive असल्याने search पुढच्या candidate ला जात राहते.
सर्वांत डाव्या redex चे न्यूनीकरण या बाह्य search ला अखंड
उलगडते; बाह्य function position मधील search संपत नाही.
त्याला सामान्य रूप असते, तर normalizing leftmost strategy
ते सांत पायऱ्यांत शोधली असती. यासाठी
\href{https://www.cs.ox.ac.uk/andrew.ker/docs/lambdacalculus-lecture-notes-ht2009.pdf}{Andrew D. Ker, \emph{Lambda Calculus and Types}, प्रमेय 3.2.3,
मुद्रित पृष्ठ 35} मधील standardization परिणाम वापरला आहे.
या search ला बाहेरून अधिक पदे उपयोजली तरी ते बाह्य
function position मध्येच राहते; अनंत search टाकून दिले जात नाही.
आता आधीच्या \ref{lam:int:lrp:thm:computable-lambda} ची सिद्धता
पूर्ण करू. क्लिनी प्रमाण रूपानुसार
\[ f(\vec x)\simeq e(\umin{n}{r(\vec x,n)}),
\]
येथे $e$ आणि $r$ सर्वत्र परिभाषित आदिम पुनरावर्ती फलने
आहेत; zero test मान्य halting record ओळखते. त्यांची बंद
प्रतिनिधी पदे $E$ आणि $R$ आधीच्या total constructions ने
मिळतात. वरील search construction मध्ये अनेक निविष्ट प्राचले
तशीच राखून त्याच्या लघुतमीकरणाचा प्रतिनिधी $W$ मिळतो.
मग $I=\lambd[t][t]$ घेऊन
\[ F(\vec x)=(\lambd[v][v\,I\,(E v)])\,W(\vec x)
\]
असे घ्या; आवश्यक बद्ध चरे ताजी निवडली आहेत. Search
ने $s$ हा अंक दिल्यास
\[ (\lambd[v][v\,I\,(E v)])\,\num{s}
 \red \num{s}\,I\,(E\,\num{s})
 \red E\,\num{s}.
\]
कारण ओळख फलनाचे कितीही पुनरावर्तन केले तरी value बदलत
नाही. शेवटचे पद योग्य output अंक देते. Search अखंड चालत
असेल, तर ते बाह्य function position मध्ये force होत राहते;
ते टाकून देणाऱ्या extractor पर्यंत पोहोचत नाही आणि सामान्य
रूप नसते. याने \ref{lam:int:lrp:thm:computable-lambda} स्वतंत्रपणे
सिद्ध झाले.
शेवटी प्रारंभीच्या उपपत्तीत $f$ आंशिक असण्याची सामान्य
शक्यता घ्या. \ref{lam:int:cmp:thm:lambda-computable} ने ते आंशिक
संगणनक्षम आहे. Strict अपरिबद्ध लघुतमीकरणातील प्रत्येक
पूर्वीचे value परिभाषित व positive असावे लागते; हा शोध
आंशिक संगणनक्षम आहे. पूर्ण झालेल्या
\ref{lam:int:lrp:thm:computable-lambda} ने त्याला बंद लॅम्डा
प्रतिनिधी मिळतो. Arbitrary आंशिक प्रतिनिधीला थेट वरील
selector मध्ये बसवणे पुरेसे आहे असे गृहीत घेतलेले नाही.
\end{proof}



